\section{Discussion and limitations}

Our results bear on how evidence about digital minds should be collected. A
self-report is the cheapest signal available. It is also the signal most readily
read as a window into a model's states. In our probe the report and
the state came apart under an intervention the mind could not detect. Any
evaluation that treats a model's account of its own reasoning as evidence about
that reasoning therefore needs a way to check the account against something else.
Building the mind out of parts gives that second source, because the experimenter
holds the window the report was written from.

The framing also changes what a preference is attached to. When a mind is one
model call, a stated preference belongs to the model. When a mind is a subject, a
stage, and an ego, the same stated preference may belong to the part that spoke,
to the window it was handed, or to the composition. That distinction matters for
the question of which entity is the object of moral concern, and it is
unavailable to a design that treats the system as a single box. We therefore
conjecture that architectural probes are the practical route to separating a
model's dispositions from the character its context invites it to play.

\subsection*{Limitations}

The pilot runs on one small open-weight model, so we cannot say whether the
dissociation we observe survives at frontier scale. Larger models may report
their context more faithfully, and the effect may reverse. The probe measures a
planted commitment by string match, which is objective but narrow. It detects
whether the mind names the planted content and not whether it describes the same
content in other words, so our misattribution count is a lower bound on how often
the report tracks the window.

The design also assumes that a pipeline of separate model calls is a reasonable
stand-in for a mind with parts. We state that assumption because our
reading of the result depends on it. If the architecture we impose bears no
relationship to whatever structure the underlying model has, then our probe
describes the composition we built and not the model inside it. We read our result in the weaker way, as a statement
about what a composed system reports, and we think the stronger reading needs
evidence that the parts correspond to something.

Two further gaps are worth naming. We study one architecture and one kind of
intervention, so we do not know how the effect varies with the shape of the mind.
We also do not address welfare or valence at all, which is a large part of what
the sprint asks about. Our contribution is an instrument and one demonstration of
what it can separate.

\subsection*{Future work}

The natural next step is scale. Running the same probe across a ladder of model
sizes would show whether reports become more faithful as capability grows, which
is the question an evaluator actually needs answered. A second step is to vary the
intervention. Removing reasoning, contradicting it, and attributing it to someone
else are different manipulations, and a mind that reports faithfully under one may
not under the others. A third step is to give the mind a workspace that broadcasts
to its parts, which turns an architectural indicator drawn from theories of
consciousness into something a run can be checked against. Each of these is an
experiment the runtime already supports.
